% Week naming rule:
% 2026-06-w01 means the Monday of this report week is in June 2026,
% and it is the 1st Monday in that month.
% Date range: Mon 2026-06-01 to Sun 2026-06-07.

\section{June 2026 Week 1}

\begin{frame}{Reporting Period: June 2026 Week 1}
  \begin{center}
    \small Date range: Mon 2026-06-01 to Sun 2026-06-07 \\
    \scriptsize Weekly meeting logic: why this matters $\rightarrow$ what we are doing $\rightarrow$ goal $\rightarrow$ progress. \\
    \scriptsize Updates compiled on 2026-06-04/05 belong to this same Week 1 report.
  \end{center}
\end{frame}

\begin{frame}{Recall: The Two Core Problems}
  \small
  \begin{block}{Problem 1: long-context inference}
    At test time, the model may need to answer from very long evidence: logs,
    traces, documents, retrieved cases, or code. The question is how to use that
    evidence without always paying full-context cost.
  \end{block}
  \begin{block}{Problem 2: catastrophic forgetting during training}
    When we train or adapt a model for a specialized setting, it may improve on
    that setting but lose general abilities such as reasoning, retrieval, or
    instruction following.
  \end{block}
  \begin{itemize}
    \item Root-cause-analysis (RCA) benchmarks matter here because they naturally
      combine both problems: long evidence at inference time and specialized
      adaptation pressure at training time.
  \end{itemize}
\end{frame}

\begin{frame}{Preliminaries: Vocabulary For This Talk}
  \scriptsize
  \begin{tabular}{p{0.28\textwidth}p{0.64\textwidth}}
    \toprule
    Term & Meaning in this project \\
    \midrule
    Frozen base model & The original LLM is not updated; this protects its general ability. \\
    Adaptation module & A small extra component trained for the task while the base stays frozen. \\
    Learned memory module / wrapper & The adaptation module tested here; it compresses long evidence into soft memory tokens. \\
    Soft memory / soft prompt & Continuous vectors inserted into the model, not human-readable text. \\
    Gate & A decision rule that chooses whether to use the learned memory or fall back to the base/full-context path. \\
    Do-no-harm & The gate should keep gains when memory helps and avoid losing points when memory is unsafe. \\
    \bottomrule
  \end{tabular}
\end{frame}

\begin{frame}{Background: Why Study These Two Problems?}
  \begin{block}{Problem}
    RCA-style tasks are a motivating benchmark family because they expose both
    long-context inference and training-time forgetting. They require long
    evidence at test time, and task-specific training can still damage general
    model behavior.
  \end{block}
  \begin{itemize}
    \item We want task adaptation without overwriting the frozen base model.
    \item The useful module should compress long context, be reusable, and be
      safe to turn off when it is not appropriate.
    \item This connects May Week 4 to June Week 1: first build a wrapper, then
      measure when it is safe to use.
  \end{itemize}
\end{frame}

\begin{frame}{What Is The Use?}
  \begin{columns}[T]
    \begin{column}{0.49\textwidth}
      \begin{block}{Practical use}
        A learned memory module could make long-context inference cheaper and more modular:
        keep the base model frozen, attach a small learned memory path, and reuse
        compressed evidence.
      \end{block}
      \begin{itemize}
        \item Avoid full-context cost when gist-level memory is enough.
        \item Preserve a fallback path to the original base model.
      \end{itemize}
    \end{column}
    \begin{column}{0.47\textwidth}
      \begin{block}{Research use}
        It gives a clean testbed for frozen-base adaptation: what can a fixed-size
        soft memory encode, and when does it fail?
      \end{block}
      \begin{itemize}
        \item Positive result: wrapper can learn.
        \item Boundary result: fixed memory has a capacity wall.
      \end{itemize}
    \end{column}
  \end{columns}
\end{frame}

\begin{frame}{What Are We Doing?}
  \small
  \begin{enumerate}
    \item \textbf{Train a learned memory module} around a frozen base model: long evidence
      context $c_{1:n}$ is compressed into soft memory $m$.
    \item \textbf{Evaluate the wrapper} against no-context, full-context,
      retrieval, and Gist-style baselines.
    \item \textbf{Characterize the boundary}: where memory helps, where it is
      neutral, and where it collapses.
    \item \textbf{Add a gate}: only let memory participate when intrinsic signals
      suggest compression is safe.
  \end{enumerate}
  \vspace{0.5em}
  \begin{block}{Short version}
    We are not claiming ``memory replaces context.'' We are building controlled
    frozen-base adaptation: learn memory, measure its limits, then gate it.
  \end{block}
\end{frame}

\begin{frame}{Goal For This Line Of Work}
  \begin{block}{Main goal}
    A frozen-base memory system that improves long-context cases where
    compression is useful, while preserving the base model when memory is unsafe.
  \end{block}
  \begin{itemize}
    \item \textbf{v1}: characterize the learned soft-prompt wrapper; write the
      result as a bit-capacity-wall paper.
    \item \textbf{v1.5}: build a do-no-harm gate between wrapper and base.
    \item \textbf{v2}: extend from one-shot compressed memory to cross-session
      latent memory with read/write behavior.
  \end{itemize}
\end{frame}

\begin{frame}{Architecture: Input, Memory, Gate, Output}
  \centering
  \small
  \begin{tikzpicture}[
    scale=0.82,
    transform shape,
    node distance=0.75cm and 0.75cm,
    box/.style={draw, rounded corners, align=center, minimum width=2.3cm, minimum height=0.85cm, fill=blue!7},
    train/.style={draw, rounded corners, align=center, minimum width=2.4cm, minimum height=0.85cm, fill=green!9},
    warn/.style={draw, rounded corners, align=center, minimum width=2.4cm, minimum height=0.85cm, fill=red!8},
    gate/.style={draw, diamond, aspect=2.1, align=center, inner sep=1.2pt, fill=orange!16},
    arr/.style={-{Latex[length=2mm]}, thick},
    dashedarr/.style={-{Latex[length=2mm]}, thick, dashed}
  ]
    \node[box] (ctx) {Long evidence\\context $c_{1:n}$};
    \node[box, below=of ctx] (query) {Query $q$};
    \node[train, right=0.85cm of ctx] (wrapper) {Trainable wrapper\\$G_\phi$};
    \node[box, right=0.85cm of wrapper] (mem) {Soft memory\\$m$, $K=32$};
    \node[gate, right=0.82cm of mem] (gate) {Gate\\open?};
    \node[box, above right=0.55cm and 0.82cm of gate] (basew) {Frozen base\\$F_\theta(q,m)$};
    \node[box, below right=0.55cm and 0.82cm of gate] (base0) {Frozen base\\$F_\theta(q)$ or\\full context};
    \node[box, right=0.82cm of $(basew)!0.5!(base0)$] (out) {Answer $y$\\+ decision trace};

    \draw[arr] (ctx) -- node[above]{enc.} (wrapper);
    \draw[arr] (wrapper) -- (mem);
    \draw[arr] (mem) -- (gate);
    \draw[arr] (query.east) -| (gate.south);
    \draw[arr] (gate) -- node[above, sloped]{safe} (basew);
    \draw[arr] (gate) -- node[below, sloped]{unsafe} (base0);
    \draw[arr] (basew) -- (out);
    \draw[arr] (base0) -- (out);

    \node[warn, below=1.0cm of mem] (risk) {Risk: forced memory\\can hurt exact/unseen tasks};
    \draw[dashedarr, red!70!black] (risk) -- (gate);
  \end{tikzpicture}
\end{frame}

\begin{frame}{How Is It Going? v1 Result}
  \small
  \begin{table}
    \centering
    \begin{tabular}{lccc}
      \toprule
      Benchmark & Wrapper & Gist & Best reference \\
      \midrule
      QuALITY & 0.193 $\pm$ 0.032 & 0.180 $\pm$ 0.044 & no-context 0.141 \\
      MuSR-mm & 0.493 $\pm$ 0.008 & 0.501 $\pm$ 0.013 & full-context 0.551 \\
      RULER-NIAH & 0.000 $\pm$ 0.000 & 0.000 $\pm$ 0.000 & full-context 0.995 \\
      \bottomrule
    \end{tabular}
  \end{table}
  \begin{itemize}
    \item \textbf{Good news}: the wrapper can train and helps on compressible
      gist-style inputs.
    \item \textbf{Lesson}: it is not a universal replacement for long context;
      exact-retrieval cases expose a capacity wall.
    \item \textbf{Setting}: P08-S2; full facts in
      \href{https://github.com/BDeMo/research-notes/blob/main/notes/plans/08-model-outputs-delta-w/v1-results-2026-06-03.md}{v1 results summary}.
  \end{itemize}
\end{frame}

\begin{frame}{How Is It Going? Interpretation}
  \begin{enumerate}
    \item \textbf{Wrapper can learn}: memory is useful when the answer can be
      represented as compressed/gist information.
    \item \textbf{Wrapper changes inference}: it is not neutral; the learned
      memory path can pull the base away from safer behavior.
    \item \textbf{Always-on memory is risky}: if the task needs exact details or
      unseen reasoning transfer, forced memory can lose information.
    \item \textbf{Next design}: put a gate between wrapper and base, and open it
      only when intrinsic signals say compression is safe.
  \end{enumerate}
\end{frame}

\begin{frame}{How Is It Going? v1.5 Gate}
  \small
  \begin{table}
    \centering
    \begin{tabular}{llcc}
      \toprule
      Signal family & Interpretation & Direction & Use \\
      \midrule
      Confidence / seq logprob & wrapper likely correct & positive & open \\
      Wrapper-to-base KL / JS / TV & wrapper fights base & inverse & close \\
      Memory/write dynamics & wrapper competence signal & positive & open/close \\
      Mid/late-layer drift & late reasoning distorted & inverse & close \\
      \bottomrule
    \end{tabular}
  \end{table}
  \begin{itemize}
    \item \textbf{Goal}: keep in-distribution gains while avoiding OOD harm.
    \item \textbf{Setting}: P08-S3/P08-S6 for signal probes; detailed v1.5
      results continue in the appended gate section.
    \item Main deliverables:
      \href{https://github.com/BDeMo/research-notes/tree/main/notes/plans/08-model-outputs-delta-w/raw/grids-2026-06-04}{signal grids} and
      \href{https://github.com/BDeMo/research-notes/blob/main/notes/plans/08-model-outputs-delta-w/summary/2026-06-04/v1.5-intrinsic-gating-study-2026-06-04.pdf}{v1.5 PDF}.
  \end{itemize}
\end{frame}

\begin{frame}{This Week's Status Summary}
  \begin{columns}[T]
    \begin{column}{0.48\textwidth}
      \begin{block}{Done}
        \begin{itemize}
          \item v1 reframed as bit-capacity characterization.
          \item Main result cells, settings, and terminology are linked.
          \item Gate direction is now concrete, not vague.
        \end{itemize}
      \end{block}
    \end{column}
    \begin{column}{0.48\textwidth}
      \begin{block}{Next}
        \begin{itemize}
          \item Polish v1 paper story.
          \item Run/validate do-no-harm gate.
          \item Test if multi-task training widens the safe region.
        \end{itemize}
      \end{block}
    \end{column}
  \end{columns}
  \vspace{0.6em}
  \begin{block}{Meeting takeaway}
    The project is moving from ``can memory learn?'' to ``can memory be used
    safely?'' The current answer is positive but bounded: useful compression
    exists, and the gate is the mechanism that can make it deployable.
  \end{block}
\end{frame}

\begin{frame}{Where To Read More}
  \small
  \begin{tabular}{p{0.30\textwidth}p{0.62\textwidth}}
    \toprule
    Item & Link \\
    \midrule
    Plan 08 hub &
    \href{https://github.com/BDeMo/research-notes/blob/main/notes/plans/08-model-outputs-delta-w/README.md}{README / deliverables hub} \\
    v1 facts &
    \href{https://github.com/BDeMo/research-notes/blob/main/notes/plans/08-model-outputs-delta-w/v1-results-2026-06-03.md}{main results summary} (setting P08-S2) \\
    v1.5 report &
    \href{https://github.com/BDeMo/research-notes/blob/main/notes/plans/08-model-outputs-delta-w/summary/2026-06-04/v1.5-intrinsic-gating-study-2026-06-04.pdf}{compiled PDF} (setting P08-S3) \\
    v1.5 signal grid &
    \href{https://github.com/BDeMo/research-notes/tree/main/notes/plans/08-model-outputs-delta-w/raw/grids-2026-06-04}{CSV + figures} (setting P08-S3) \\
    v2 setting &
    \href{https://github.com/BDeMo/research-notes/blob/main/notes/plans/08-model-outputs-delta-w/v2-plan.md}{v2 plan} (setting P08-S4) \\
    \bottomrule
  \end{tabular}
\end{frame}
